\section{总结与延伸}

\subsection{讲者的核心总结}

\begin{figure}[H]
\centering
\includegraphics[width=0.95\textwidth]{slides-images/slide-074.jpg}
\caption{Lecture 14 总结：推理和智能体两大主题的关键要点\protect\footnotemark}
\end{figure}
\footnotetext{Slides 第75页。}

Shikhar Murty 在课程结尾将两大主题的要点归纳如下：

\textbf{推理（Reasoning）：}
\begin{itemize}
    \item 可以通过\textbf{提示}（CoT、Self-Consistency、Least-to-Most）激发推理行为
    \item 可以通过\textbf{蒸馏}（Orca）将大模型的推理能力迁移到小模型
    \item 可以通过\textbf{自我训练}（ReST$^{EM}$）让模型在自己的输出上迭代改进
    \item \textbf{反事实评估}表明推理可能并不系统——记忆和模式匹配起了很大作用
\end{itemize}

\textbf{智能体（Agents）：}
\begin{itemize}
    \item 主要使用\textbf{提示和上下文学习}
    \item \textbf{BAGEL} 通过探索和迭代重标注生成合成训练数据
    \item \textbf{多模态}模型（LLaVA、Pix2Struct）使直接处理视觉输入成为可能
    \item 当前\textbf{评估基准仍极具挑战性}，模型表现与人类存在巨大差距
\end{itemize}

\subsection{全课知识图谱}

\begin{center}
\begin{tikzpicture}[
    node distance=0.9cm,
    every node/.style={draw, rounded corners, minimum width=3.5cm, minimum height=0.7cm, font=\small},
    arrow/.style={-{Stealth[length=3mm]}, thick}
]
% Reasoning branch
\node (reason) {\textbf{推理（Reasoning）}};
\node (prompt) [below left=1cm and 0.5cm of reason] {提示方法};
\node (train) [below=of reason] {训练方法};
\node (eval) [below right=1cm and 0.5cm of reason] {评估与反思};

% Agent branch
\node (agent) [below=2.5cm of reason] {\textbf{智能体（Agents）}};
\node (pre) [below left=1cm and 0.5cm of agent] {Pre-LLM 方法};
\node (llm) [below=of agent] {LLM 智能体};
\node (data) [below right=1cm and 0.5cm of agent] {合成数据};

\draw[arrow] (reason) -- (prompt);
\draw[arrow] (reason) -- (train);
\draw[arrow] (reason) -- (eval);
\draw[arrow] (reason) -- (agent);
\draw[arrow] (agent) -- (pre);
\draw[arrow] (agent) -- (llm);
\draw[arrow] (agent) -- (data);

\node[draw=none, right=0.3cm of prompt, text width=4cm, font=\scriptsize] {CoT, Self-Consistency, Least-to-Most};
\node[draw=none, left=0.3cm of eval, text width=4cm, font=\scriptsize, align=right] {忠实性测试, 反事实评估};
\node[draw=none, right=0.3cm of pre, text width=4cm, font=\scriptsize] {语义解析, 计划推断, RL};
\node[draw=none, left=0.3cm of data, text width=4cm, font=\scriptsize, align=right] {BAGEL: 探索+重标注};
\end{tikzpicture}
\end{center}

\subsection{两大主题的深层联系}

\begin{importantbox}{推理与智能体的统一视角}
本讲的两个看似独立的主题实际上有深层联系：
\begin{enumerate}
    \item \textbf{技术共享}：推理中的自我训练（ReST$^{EM}$）和智能体中的合成数据生成（BAGEL）使用了相同的核心思路——生成、过滤、迭代改进
    \item \textbf{能力依赖}：智能体需要推理能力来进行规划和决策；推理能力的缺陷直接限制了智能体的表现
    \item \textbf{共同挑战}：两者都面临\textbf{系统性}的问题——推理的``记忆vs推理''困境对应着智能体的``提示鸿沟''
    \item \textbf{架构统一}：两者都可以被建模为自回归序列生成——推理是生成推理链，智能体是生成动作轨迹
\end{enumerate}
\end{importantbox}

\subsection{关键 Takeaways}

\begin{importantbox}{五条核心原则}
\begin{enumerate}
    \item \textbf{中间步骤至关重要}：无论是推理链还是动作轨迹，让模型``展示工作过程''都能显著提升性能
    \item \textbf{自我改进是可能的}：模型可以通过在自己的输出上训练来超越人类标注的上限
    \item \textbf{推理可能不系统}：反事实评估揭示了 LLM ``推理''中记忆和模式匹配的成分
    \item \textbf{决策即语言建模}：智能体的动作预测本质上就是条件语言生成
    \item \textbf{差距仍然巨大}：在推理和智能体两个领域，模型与人类之间的差距都还很大，还有大量改进空间
\end{enumerate}
\end{importantbox}

\subsection{拓展阅读}

\begin{itemize}
    \item Wei et al., Chain-of-Thought Prompting: \url{https://arxiv.org/abs/2201.11903}
    \item Wang et al., Self-Consistency: \url{https://arxiv.org/abs/2203.11171}
    \item Zhou et al., Least-to-Most Prompting: \url{https://arxiv.org/abs/2205.10625}
    \item Kojima et al., Zero-shot CoT: \url{https://arxiv.org/abs/2205.11916}
    \item Mukherjee et al., Orca: \url{https://arxiv.org/abs/2306.02707}
    \item Singh et al., ReST$^{EM}$: \url{https://arxiv.org/abs/2312.06585}
    \item Wu et al., Reasoning or Reciting: \url{https://arxiv.org/abs/2307.13702}
    \item Hodel et al., Analogical Reasoning Counterfactuals: \url{https://arxiv.org/abs/2402.02442}
    \item Chen et al., Decision Transformer: \url{https://arxiv.org/abs/2106.01345}
    \item Zhou et al., WebArena: \url{https://arxiv.org/abs/2307.13854}
    \item Liu et al., LLaVA: \url{https://arxiv.org/abs/2304.08485}
    \item Lee et al., Pix2Struct: \url{https://arxiv.org/abs/2210.03347}
    \item Murty et al., BAGEL: \url{https://arxiv.org/abs/2401.15839}
\end{itemize}

\end{document}
